\documentclass[10pt,a4paper]{amsart}
\usepackage[margin=19mm]{geometry}
\usepackage{amsmath,amssymb,mathtools}
\usepackage[scr=rsfs,cal=euler]{mathalfa}
\usepackage{xfrac}
\usepackage[unicode,hidelinks,bookmarks=false]{hyperref}
\newcommand{\ie}{i.e.\ }
\newcommand{\cf}{cf.\ }
\newcommand{\resp}{respectively\ }
\newcommand{\Subsection}{\subsection}
\providecommand{\nobreakdash}{\nobreak\hbox{-}}
\setlength{\emergencystretch}{2em}
\multlinegap0pt
\usepackage{bm}
\usepackage{bbm}
\usepackage{dsfont}
\usepackage{xspace}
\usepackage{cancel}
\usepackage{mathtools}
\usepackage{amsthm}
\makeatletter
\@ifundefined{claim}{\newtheorem{claim}{Claim}}{}
\@ifundefined{corollary}{\newtheorem{corollary}{Corollary}}{}
\makeatother
\makeatletter
\expandafter\ifx\csname proof\endcsname\relax
\newenvironment{proof}[1][Proof]{\textbf{#1.} }{\ \rule{0.5em}{0.5em}}
\fi
\makeatother
\title[Advances on the Packing Coloring Conjectures of Subcubic Gra]{Advances on the Packing Coloring Conjectures of Subcubic Graphs}
\author{Ayman El Zein, Maidoun Mortada}
\date{Source: arXiv:2503.20239v3 (CC BY 4.0)}
\begin{document}
\maketitle

\noindent\emph{Source and licence.} Advances on the Packing Coloring Conjectures of Subcubic Graphs. Version arXiv:2503.20239v3 (CC BY 4.0), distributed under the Creative Commons Attribution 4.0 International licence, as recorded in the arXiv licence metadata for this exact version and shown on the abstract page https://arxiv.org/abs/2503.20239v3. Full rights notice: licenses/arxiv-2503.20239-CC-BY-4.0.txt.\par\smallskip

\noindent\textit{Editorial scope.} Counted unit: the Corollary whose source label is \texttt{None} in the article (source0.tex lines 427--430, source label \texttt{None}), delivered together with the article's own complete original proof of it (source0.tex lines 431--458, one \texttt{proof} environment of 28 lines). No mathematics is written, repaired or re-derived: statement and proof are the article's own text line for line. Required context retained uncounted: none. Every definition, display and label of those lines is retained unchanged. Omitted from this package and not redistributed: 1--426, 459--545. Nothing inside a retained excerpt is omitted or altered: every retained body is delivered exactly as the article's lines are written, so no omission receipt of any kind accompanies this package. Citations: the retained excerpts carry no citation command at all, so no bibliography entry of the article is transcribed and no reference is redistributed. Reference adaptation: none was needed and none was made; every reference inside the delivered unit resolves inside it, so no unresolved reference or citation is produced. Printed slips kept literal: no printed word, symbol, hypothesis or number of the counted statement or of its proof was altered. Layout and class: the article's own class is \texttt{scrartcl} with packages including babel, fontenc, amssymb, amsmath, hhline, longtable, amscd, theorem, array, delarray, multicol, makeidx, float, tikz; the wrapper is the lane's pinned \texttt{amsart} editorial wrapper at 10pt on A4, which restates the theorem-like environments the retained text needs and adds the article's own macro definitions that the retained text needs (none), with extra wrapper packages (bm, bbm, dsfont, xspace, cancel, mathtools). The delivered numbering is the unit's own. Assets: no retained span contains a figure environment, a graphics-inclusion command, a PDF-page inclusion, a table or a TikZ picture; no image file is copied into this package, the package declares no asset dependency and no image file is redistributed with it. Licence: version arXiv:2503.20239v3 of this article is distributed under the Creative Commons Attribution 4.0 International licence, as recorded in the arXiv licence metadata for this exact version and shown on the abstract page https://arxiv.org/abs/2503.20239v3. A full rights notice is delivered at licenses/arxiv-2503.20239-CC-BY-4.0.txt. Characters: every retained excerpt is pure ASCII, so the delivered engine, XeLaTeX with its default Latin Modern text font, reports no missing character.
\begin{corollary}
Let $G$ be a connected cubic graph that contains a bridge $e = xy$ such that each of $x$ and $y$ belongs to a cycle of length at most four. 
Then, $G$ admits a $(1,1,2,2)$-packing coloring.
\end{corollary}

\begin{proof}
Let $G' = G - e$. By Theorem~1, $G'$ admits a $(1,1,2,2)$-packing coloring. Fix such a coloring $c : V(G') \to \{1_a,1_b,2_a,2_b\}$. If $c$ is still valid after restoring $e$, we are done. Otherwise, the only possible conflicts in $G$ created by restoring $xy$ are:
\begin{enumerate}
  \item[\textup{(A)}] $x$ and $y$ receive the same color;
  \item[\textup{(B)}] $x$ and a neighbor of $y$ have the same $2$-color;
  \item[\textup{(C)}] $y$ and a neighbor of $x$ have the same $2$-color.
\end{enumerate}
 Let $y_1$ and $y_2$ be the neighbors of $y$ other than $x$. Here is the treatment of each conflict when it appears:\begin{itemize}
 \item \textbf{Resolution of conflict (A).}
 In this case, it is enough to switch the colors $1_a$ and $1_b$ (resp. $2_a$ and $2_b$) if $x$ and $y$ have the same 1-color (resp. 2-color) in the connected component containing $x$ in $G'$. Hence, conflict (A) is always resolvable.
 
\item \textbf{Resolution of conflict (B) when conflict (C) doesn't occur.} 
 Suppose $x$ and $y_1$ have the same 2-color, say $2_a$. If $y_2$ is not colored by $2_b$, then switch $2_a$ and $2_b$ in the  component of $x$ in $G'$ and so the conflict is removed. Else; that is $y_2$ is of color $2_b$. In this case, a global switch in the component of $x$ is not enough to resolve the conflict. Because $y$ lies on a cycle $C$ of length $3$ or $4$, we will show that we can locally recolor to make one of $y_1,y_2$ a $1$-colored vertex and then resolve the conflict accordingly.

If $G[\{y,y_1,y_2\}]$ is a triangle, then recolor $y_1$ by $1_i$, $i\in\{a,b\}$, such that $y_i$ has no neighbor of color $1_i$ in $G\setminus \{y\}$ and recolor $y$ by $1_j$ such that $j\in \{a,b\}\setminus\{i\}$. This removes the conflict. 

For the case when  $G[\{y,y_1,y_2\}]$ is not a triangle, then $y$, $y_1$ and $y_2$ are contained in a  square. Let $z$ be the fourth vertex of this square. If there exists $i\in\{a,b\}$ such that none of the neighbors of $y_1$ in $G\setminus \{y\}$ is of color $1_i$ then proceed as in the previous case; that is  recolor $y_1$ by $1_i$ and recolor $y$ by $1_j$ such that $j\in \{a,b\}\setminus\{i\}$. Else, if there exists $i\in\{a,b\}$ such that none of the neighbors of $y_2$ in $G\setminus \{y\}$ is of color $1_i$, then recolor $y_2$ by $1_i$,   recolor $y$ by $1_j$ such that $j\in \{a,b\}\setminus\{1_i\}$, and switch the colors $2_a$ and $2_b$ in the  component of $x$ in $G'$. Else, both colors $1_a$ and $1_b$ are given to the neighbors of $y_1$ (resp. $y_2$) in $G\setminus \{y\}$. Let $z_1$ (resp. $z_2$) be the third neighbor of $y_1$ (resp. $y_2$). Hence, $z$ is of color $1_l$ and $z_1$ and $z_2$ are both of color $1_k$ such that $\{1_l,1_k\}=\{1_a,1_b\}$. We have the following observation:
\begin{claim}
    The vertex $z$ can be given a color other than $1_l$ and then $y_1$ or $y_2$ can be recolored in order to remove the conflict.
\end{claim}
\begin{proof}
If $z$ has no neighbor of color $1_k$,  recolor $z$ by $1_k$.  Then,  none of the neighbors of $y_1$ in $G\setminus \{y\}$ is now of color $1_l$ and so the conflict can be removed as before by recoloring $y_1$ by $1_l$ and $y$ by $1_k$. Else, let $z'$ be the neighbor of $z$ of color $1_k$. If $z'$ has no neighbor other than $z$ of color $1_l$, then recolor $z'$ by $1_l$, $z$ by $1_k$, $y_1$ by $1_l$, and $y$ by $1_k$. Hence, the conflict is removed.  Otherwise, $z'$ has two neighbors of color $1_l$ and so, since   $G[\{y,y_1,y_2,z\}]$ is a square, then there exists $i\in \{a,b\}$ such that $z$ is at distance at least three from each vertex of color $2_i$ in $G\setminus \{y_1,y_2\}$. If $i=a$, recolor $z$ by $2_a$, $y_1$ by $1_l$, $y$ by $1_k$, and so the conflict is removed. Else; that is $i=b$, then recolor $z$ by $2_b$, $y_2$  by $1_l$, $y$ by $1_k$, and finally switch the two colors $2_a$ and $2_b$ in the connected component containing $x$. 
\end{proof}


\item \textbf{Resolution when both (B) and (C) occur.}
In this case, there exists $i\in \{a,b\}$ (resp. $j\in\{a,b\}$) such that $x$ (resp $y$) has no neighbor of color $1_i$ (resp $1_j$). Thus, recolor $x$ by $1_i$. For $y$, recolor $y$ by $1_j$ if $i\neq j$. Else; that is $i=j$, switch $1_a$ and $1_b$ in the connected component containing $y$ in $G'$ and then recolor $y$ by $1_l$ such that $l\neq i$.\end{itemize}
\end{proof}

\end{document}
